\section{Conclusion}
\label{sec:conclusion}

This study holds the FedAvg AE ($E{=}1$, full participation), training protocol, and seeds fixed across B1--B4;
only the post-training threshold calibration scope varies.
On heterogeneous N-BaIoT devices, per-client calibration (B2) provides the strongest
FPR-dispersion reduction ($\Delta = 0.718$, all five seed deltas positive, range $[0.588,\,0.772]$);
cluster-mean calibration (B4) recovers a meaningful but partial share
($\approx$52\% of B2's improvement) without requiring a device taxonomy, at the cost
of transmitting distributional metadata that is not private.
Family-mean calibration (B3) performs poorly: coarse device-type labels do not align
with reconstruction-error calibration structure.

When benign error distributions are near-homogeneous (CICIoT2023 under the file-level
pseudo-client partition, Dirichlet near-IID), no dispersion reduction is observed---this is
an applicability boundary rather than evidence that personalization is unnecessary for these datasets generally.
The main tradeoff under heterogeneous conditions is that B2 reduces FPR dispersion
and improves mean Macro-F1, but the lower-tail P10 Macro-F1 degradation is concentrated
in the low-separability Ennio Doorbell client---a client-specific failure of percentile
calibration rather than evidence of a universal FPR--TPR tradeoff.
Per-client calibration is an FPR-equity intervention, not a uniform detection-quality improvement.

\textbf{Practitioner guidance.}
First, measure benign calibration-error heterogeneity across clients, for example
using pairwise distribution divergence or dispersion of local calibration thresholds.
Second, choose the threshold policy according to that heterogeneity and deployment
constraints: use B1 when clients are near-homogeneous or per-client state is
undesirable; use B2 when local calibration is available and false-alarm equity is the
priority; use B4 when taxonomy-free grouping is acceptable, while accounting for the
metadata leakage introduced by shared distributional fingerprints.

\textbf{Limitations.}
Regime~A uses only 9 physical devices (B4 is exploratory at this scale);
Regime~B uses randomized file-level pseudo-clients (not physical devices);
$E{=}1$ with full participation limits local drift---results characterize threshold personalization under a frequently synchronized encoder, not drift-heavy settings;
five seeds limit statistical power; no adversarial robustness, formal privacy, or hardware evaluation is provided.
Evaluating $E > 1$ and alternative CICIoT2023 partitions (physical-device, attack-source) is future work.
Threshold personalization is complementary to---not a replacement for---model-level personalization.

\textbf{Future work:} larger physical fleets, $E > 1$ sensitivity,
privacy-preserving grouping, model-level personalization, adversarial robustness,
and deployment validation under concept drift.

\noindent\textbf{Availability.}
A replication package containing source code, configuration files, per-seed results,
and table/figure generation scripts is archived on Zenodo (DOI: 10.5281/zenodo.20315654).
Detailed reproduction commands are provided in the repository documentation.
